In log-ratio units, the unweighted slopes per unit of wage parity are $-2.05$ for $P$ ($t=-2.5$) and $+0.84$ for $C$ ($t=+1.4$). These components use full top-decile totals, including Chile, as in Table \ref{tab:crossfield_decomposition}. Figure \ref{fig:gender_norms} instead plots female-minus-male percentage-point gaps and conditions Chile's application rates on submitted lists. Pipeline precision is similar on either scale ($|t|=2.54$ in log points versus $2.54$ in percentage points). Leave-one-setting-out slopes on $C$ range from $-0.85$ to $+1.02$. Excluding Uganda changes the slope on $C$ to $-0.85$ ($t=-0.6$); both the log-ratio and percentage-point choice slopes change sign. The persistent choice gap is robust across settings; the direction and precision of its association with wage parity are sensitive to scale and sample composition.
